\section{Introduction}

Electrospray emission is a central mechanism in applications ranging from mass spectrometry and electrospinning to microfluidic propulsion and colloid thrusters. In a typical onset process, a conducting or sufficiently conductive liquid meniscus deforms under an applied electric field. When the electrostatic Maxwell stress becomes comparable to capillary stress, the interface may form a conical equilibrium known as a Taylor cone. Classical theory predicts a semi-vertical half-angle of approximately \SI{49.3}{\degree} in the ideal conducting, charge-free limit \cite{taylor1964disintegration}.

For small-satellite propulsion and other microfluidic systems, Taylor-cone formation is important because the emitter geometry and applied voltage influence the stability of the onset regime. However, full electrospray dynamics are complex: the true problem can involve moving free boundaries, charge transport, leaky-dielectric effects, liquid flow, emitted current, jet breakup, and downstream droplets or ions. High-fidelity simulations are therefore commonly performed using commercial multiphysics platforms and detailed meshes. Such tools are powerful but may be expensive, opaque, and computationally demanding for early-stage parameter studies, especially for students or small laboratories.

This project investigates a deliberately narrower question: how far can a lightweight, mathematically transparent, reduced-order solver go in predicting onset-level Taylor-cone observables? Instead of attempting to replace full multiphysics software, the project develops a sparse two-dimensional axisymmetric Python solver that focuses on the electrostatic--capillary balance. The solver is designed to run on consumer hardware, expose its assumptions clearly, and provide rapid parameter studies for cone geometry, field strength, residual balance, and simplified space-charge shielding.

The project has three validation layers. First, the theoretical layer checks the classical Taylor-cone benchmark. Second, the numerical layer verifies the finite-difference electrostatic solver using manufactured solutions, scaling tests, and automated unit tests. Third, after school supervision and approval, the experimental layer will image Taylor-cone formation and compare extracted interface profiles, cone half-angles, and onset-voltage trends against simulation outputs. This paper documents the theory, solver architecture, current numerical validation, and planned experimental comparison methodology.
